\documentclass[11pt,a4paper]{article}

\title{\textbf{Star Schema Diagram Documentation}\\[0.5cm] \textbf{AeroGrid}\\[0.3cm] Low-Altitude Airspace Intelligence Warehouse}
\author{AeroGrid Development Team}
\date{\today}

\begin{document}

\maketitle

\section{Schema Architecture}
The AeroGrid Data Warehouse utilizes a Star Schema architecture optimized for OLAP aggregations. The core structure centers around five fact tables radiating out to shared dimension tables.

\section{Core Dimensions and Star Topology}
\begin{itemize}
    \item \textbf{Star Core (\texttt{DimTime}):} Connected to all five fact tables via \texttt{time\_key}, enabling uniform temporal filtering and drill-down analysis.
    \item \textbf{Actor Star (\texttt{DimVehicle}):} Connected to tracking, journey, and close encounter facts. Uses role-playing foreign keys (\texttt{vehicle\_key\_1} and \texttt{vehicle\_key\_2}) for multi-vehicle conflict analysis in \texttt{FactCloseEncounter}.
    \item \textbf{Spatial Star (\texttt{DimSpatialGrid} and \texttt{DimAirspaceZone}):} Links position reports, close encounters, and occupancy snapshots to geographic grid cells and FAA restriction boundaries.
\end{itemize}

\section{Schema Summary}
The complete database model is formalized in DBML format (\texttt{data/diagrams/aerogrid\_schema.dbml}), ensuring seamless visualization in tools like dbdiagram.io and automated schema migration.

\end{document}
